% ======================================================================
% Macros partagées par le rapport, l'entente autonome et la présentation.
% Ce fichier est chargé (\input) dans les préambules : ne rien y composer.
% ======================================================================

% --- Noms propres et cadres ------------------------------------------------
% Utilisation : \fedesante{}, \pcap{}, \eptc{} (les accolades évitent
% que l'espace suivante soit avalée).
\newcommand{\fedesante}{FédéSanté}
% PCAP est une marque déposée du Centre de gouvernance de l'information des
% Premières Nations (CGIPN) ; le symbole ® est demandé par le CGIPN.
% TODO : vérifier la forme d'usage recommandée par le CGIPN.
\newcommand{\pcap}{PCAP\textsuperscript{\textregistered}}
% TODO : vérifier la version de l'EPTC 2 en vigueur à la date de remise
% et ajuster l'entrée correspondante dans shared/references.bib.
\newcommand{\eptc}{EPTC~2}
\newcommand{\loivingtcinq}{Loi~25}

% --- Espace réservé visible ----------------------------------------------
% Marque dans le PDF ce qui reste à rédiger. Aucun \arediger ne doit
% subsister dans la version remise (vérifier avec : grep -rn arediger).
\newcommand{\arediger}[1]{{\color{gray}\itshape[À rédiger : #1]}}

% --- Sigles et acronymes (paquet acronym) ----------------------------------
% \ac{FL} : forme longue + sigle au premier emploi, sigle ensuite.
% \acs{FL} : sigle seul ; \acl{FL} : forme longue seule.
\usepackage{acronym}
\acrodef{FL}[FL]{apprentissage fédéré (\emph{federated learning})}
\acrodef{DP}[DP]{confidentialité différentielle (\emph{differential privacy})}
\acrodef{FedAvg}[FedAvg]{moyenne fédérée (\emph{Federated Averaging})}
\acrodef{CER}[CER]{comité d'éthique de la recherche}
\acrodef{PCAP}[PCAP]{propriété, contrôle, accès et possession}
